\subsection{Regime 3: antibody in excess of target}

As in Regime 1, let $S_c=A+B$ and $S_t=A'+B$ be the total target-epitope
concentrations in the control and treatment conditions, and let
$S\in\{S_c,S_t\}$ denote either one. We assume the following for both
conditions (i.e., for each $S\in\{S_c,S_t\}$).

The assumptions are again easiest to interpret through a heuristic
estimate of the free-antibody concentration. Now there is more antibody
than target. If nearly every target epitope is bound, then about $S$
antibody molecules sit on target, and the excess $Y-S$ must go somewhere
else: either into solution ($Y_f$) or onto off-target epitopes
($Y_C+Y_D$). If off-target epitopes are far from saturated, nearly all of
them are free ($C_f+D_f\approx C+D$), so rearranging the mass-action
relation $K_{d2}=Y_f(C_f+D_f)/(Y_C+Y_D)$ gives
$$
Y_C+Y_D
=\frac{Y_f(C_f+D_f)}{K_{d2}}
\approx\frac{Y_f(C+D)}{K_{d2}}.
$$
Substituting this into ``excess $=$ free $+$ off-target-bound,''
$$
Y-S
\approx Y_f+\frac{Y_f(C+D)}{K_{d2}}
=Y_f\cdot\frac{K_{d2}+C+D}{K_{d2}},
$$
and multiplying both sides by $K_{d2}/(K_{d2}+C+D)$,
$$
Y_f\approx K_{d2}\frac{Y-S}{K_{d2}+C+D}.
$$
In words, the excess antibody is split between solution and off-target
epitopes in the ratio $K_{d2}:(C+D)$. Below, we show that this estimate
is in fact an exact lower bound on $Y_f$.

\begin{enumerate}
    \item Antibody is more abundant than target epitope:
    $$
    Y>S.
    $$
    Thus, even if every target epitope is bound, there is a positive
    excess $Y-S$ of antibody left over.

    \item Free target is scarce relative to the excess antibody:
    $$
    rS\,(K_{d2}+C+D)\ll(Y-S)^2.
    $$
    To interpret this, estimate the free target concentration
    $A_f+B_f$. Rearranging the mass-action relation
    $K_{d1}=Y_f(A_f+B_f)/(Y_A+Y_B)$, and using that about $S$ target
    epitopes are bound,
    $$
    A_f+B_f
    =K_{d1}\frac{Y_A+Y_B}{Y_f}
    \approx\frac{K_{d1}S}{Y_f}.
    $$
    Substituting the heuristic estimate of $Y_f$ and then
    $r=K_{d1}/K_{d2}$,
    $$
    A_f+B_f
    \approx\frac{K_{d1}S\,(K_{d2}+C+D)}{K_{d2}(Y-S)}
    =\frac{rS\,(K_{d2}+C+D)}{Y-S}.
    $$
    Dividing both sides of the condition by $Y-S$ shows that it is
    equivalent to
    $rS(K_{d2}+C+D)/(Y-S)\ll Y-S$: the estimated free target is much less
    than the excess antibody. This matters because every target epitope
    that is free leaves one more antibody molecule unaccounted for, so the
    antibody available to go into solution or off-target stays close to
    $Y-S$ only if the free target is small compared with $Y-S$. Usually
    this also means that the target is nearly saturated
    ($\theta_1\approx1$). The square arises because $Y-S$ appears twice:
    once in the estimate of the free target (more excess antibody drives
    more target into the bound state) and once as the yardstick against
    which the free target is compared.

    \textbf{Mirror symmetry with Regime 1.} Substituting
    $r=K_{d1}/K_{d2}$ and splitting the fraction, the left-hand side is
    $$
    rS\,(K_{d2}+C+D)
    =\frac{K_{d1}}{K_{d2}}S\,(K_{d2}+C+D)
    =K_{d1}S\left(1+\frac{C+D}{K_{d2}}\right).
    $$
    If there were no off-target epitopes ($C+D=0$), the condition would
    read $K_{d1}S\ll(Y-S)^2$. Compare Regime 1's assumption 2,
    $K_{d1}Y\ll(S-Y)^2$: the two are identical after swapping $Y$ and
    $S$. This is because the mass-action relation
    $K_{d1}=Y_f(A_f+B_f)/(Y_A+Y_B)$ treats antibody and target
    symmetrically. In Regime 1, antibody is the limiting species and target
    is in excess; in Regime 3, the roles are reversed. In both cases the
    condition says that the leftover \emph{limiting} species (free antibody
    in Regime 1, free target here) is much less than the excess of the
    \emph{abundant} species. The extra factor $1+(C+D)/K_{d2}$ breaks the
    symmetry: off-target epitopes are a sink for antibody but not for
    target, so only a fraction $K_{d2}/(K_{d2}+C+D)$ of the excess
    antibody remains free to drive target binding. A weaker push on the
    target requires a larger excess to compensate.

    Formally, below we show that $S(1-\theta_1)\le rS(K_{d2}+C+D)/(Y-S)$,
    so this condition makes the free target negligible relative to $Y-S$.

    \item Off-target epitopes are far from saturated:
    $$
    Y-S\ll K_{d2}+C+D.
    $$
    By the heuristic estimate, the off-target occupancy is
    $$
    \theta_2=\frac{Y_f}{Y_f+K_{d2}}\approx\frac{Y_f}{K_{d2}}
    \approx\frac{Y-S}{K_{d2}+C+D},
    $$
    so this condition says that the fraction of off-target epitopes that
    are bound is small. Another way to see the yardstick
    $K_{d2}+C+D$: to half-occupy the off-target epitopes, the free antibody
    must reach $K_{d2}$ and about $(C+D)/2$ antibody molecules must be bound
    off-target, for a total of about $K_{d2}+(C+D)/2$, which is within a
    factor of 2 of $K_{d2}+C+D$. The condition says that the excess
    antibody is far below this amount. Below, we show that this condition
    implies $Y_f\ll K_{d2}$.
\end{enumerate}
In summary, assumption 2 ensures that essentially all of the excess
antibody $Y-S$ is available to the ``non-target'' compartments (solution
and off-target epitopes), and assumption 3 keeps the off-target epitopes
in their unsaturated, linear-response regime. Compared with Regime 1,
there is no separate assumption limiting off-target binding: here,
off-target binding is not a perturbation to be neglected but one of the
two places where the excess antibody goes.

\textbf{Simpler sufficient condition.} Suppose $Y\ge2S$. Subtracting $S$
from both sides gives $Y-S\ge S$, and since $Y-S>0$ by assumption 1,
dividing both sides by $Y-S$ gives
$$
\frac{S}{Y-S}\le1.
$$
Writing the left-hand side of assumption 2, divided by $(Y-S)^2$, as a
product,
$$
\frac{rS\,(K_{d2}+C+D)}{(Y-S)^2}
=\frac{r(K_{d2}+C+D)}{Y-S}\cdot\frac{S}{Y-S}
\le\frac{r(K_{d2}+C+D)}{Y-S}.
$$
Therefore, if $r(K_{d2}+C+D)\ll Y-S$, then assumption 2 holds. When $Y$
approaches $S$, however, the factor $S/(Y-S)$ grows without bound, and
the simpler condition is no longer sufficient.

We now derive $\rho$ for either condition. (For the treatment condition,
$\theta_1$, $\theta_2$, $Y_f$, and $\rho$ should be read as their primed
counterparts.)

\textbf{Antibody conservation.} As in Regime 1,
$$
Y=Y_f+S\theta_1+(C+D)\theta_2.
$$

\textbf{Step 1: $\rho$ in terms of $\theta_2$.} Substituting
$K_{d1}=rK_{d2}$ into the numerator of $\rho$, and then adding and
subtracting $rY_f$,
$$\begin{aligned}
\rho
&=\frac{Y_f+K_{d1}}{Y_f+K_{d2}}
=\frac{Y_f+rK_{d2}}{Y_f+K_{d2}} \\
&=\frac{r(Y_f+K_{d2})+Y_f-rY_f}{Y_f+K_{d2}} \\
&=r+(1-r)\frac{Y_f}{Y_f+K_{d2}}.
\end{aligned}$$
Recognizing $\theta_2=Y_f/(Y_f+K_{d2})$,
\begin{equation}
\rho=r+(1-r)\,\theta_2.
\label{eq:r3-1}
\end{equation}
This is exact and holds in every regime. It says that $\rho$ interpolates
linearly between $r$ (no off-target occupancy) and 1 (full off-target
occupancy). So it remains to find $\theta_2$.

\textbf{Step 2: where the excess antibody goes.} Write the
target-bound antibody as total target minus free target,
$S\theta_1=S-S(1-\theta_1)$. Substituting into antibody conservation,
$$
Y=Y_f+S-S(1-\theta_1)+(C+D)\theta_2.
$$
Subtracting $S$ from both sides and adding $S(1-\theta_1)$ to both sides,
\begin{equation}
(Y-S)+S(1-\theta_1)=Y_f+(C+D)\theta_2.
\label{eq:r3-2}
\end{equation}
Intuitively, the antibody not on target (right-hand side) equals the
excess antibody $Y-S$ plus one antibody for each target epitope that is
left free. This is the mirror image of Regime 1's equation
\eqref{eq:r1-1}.

\textbf{Step 3: an exact expression for $\theta_2$.} From
$\theta_2=Y_f/(Y_f+K_{d2})$, multiplying both sides by $Y_f+K_{d2}$ gives
$Y_f=\theta_2(Y_f+K_{d2})$. Substituting this into the right-hand side of
\eqref{eq:r3-2} and factoring out $\theta_2$,
$$
(Y-S)+S(1-\theta_1)
=\theta_2(Y_f+K_{d2})+(C+D)\theta_2
=\theta_2\,(K_{d2}+C+D+Y_f).
$$
Dividing both sides by $K_{d2}+C+D+Y_f$,
$$
\theta_2=\frac{(Y-S)+S(1-\theta_1)}{K_{d2}+C+D+Y_f}.
$$
Factoring $Y-S$ out of the numerator and $K_{d2}+C+D$ out of the
denominator,
\begin{equation}
\theta_2
=\frac{Y-S}{K_{d2}+C+D}\cdot\frac{1+\delta_1}{1+\delta_2},
\qquad
\delta_1=\frac{S(1-\theta_1)}{Y-S},
\quad
\delta_2=\frac{Y_f}{K_{d2}+C+D}.
\label{eq:r3-3}
\end{equation}
Equations \eqref{eq:r3-1} and \eqref{eq:r3-3} together are exact. Here
$\delta_1$ is the free target relative to the excess antibody, and
$\delta_2$ is the free antibody relative to $K_{d2}+C+D$. It remains to
show that $\delta_1\ll1$ and $\delta_2\ll1$.

\textbf{Step 4: a lower bound on $Y_f$.} Because $S(1-\theta_1)\ge0$,
equation \eqref{eq:r3-2} gives
$$
Y_f+(C+D)\theta_2\ge Y-S.
$$
Dropping $Y_f$ from the denominator of $\theta_2$ makes it larger, so
$\theta_2=Y_f/(Y_f+K_{d2})\le Y_f/K_{d2}$. Substituting this on the
left-hand side (which can only make the left-hand side larger, so the
inequality still holds),
$$
Y_f+\frac{(C+D)Y_f}{K_{d2}}\ge Y-S
\quad\Longrightarrow\quad
Y_f\cdot\frac{K_{d2}+C+D}{K_{d2}}\ge Y-S.
$$
Multiplying both sides by the positive number $K_{d2}/(K_{d2}+C+D)$,
\begin{equation}
Y_f\ge K_{d2}\frac{Y-S}{K_{d2}+C+D},
\label{eq:r3-4}
\end{equation}
which is positive by assumption 1. This is the heuristic estimate from
the start of this subsection, now shown to be a lower bound.

\textbf{Step 5: $\delta_1\ll1$.} From $\theta_1=Y_f/(Y_f+K_{d1})$, as in
Regime 1, $1-\theta_1=K_{d1}/(Y_f+K_{d1})$. Dropping $K_{d1}$ from the
denominator makes the fraction larger, and then substituting the lower
bound \eqref{eq:r3-4} for $Y_f$ in the denominator makes it larger
still:
$$
1-\theta_1
=\frac{K_{d1}}{Y_f+K_{d1}}
\le\frac{K_{d1}}{Y_f}
\le\frac{K_{d1}(K_{d2}+C+D)}{K_{d2}(Y-S)}
=\frac{r(K_{d2}+C+D)}{Y-S},
$$
where the last step substitutes $r=K_{d1}/K_{d2}$. Multiplying both sides
by $S/(Y-S)$,
$$
\delta_1
=\frac{S(1-\theta_1)}{Y-S}
\le\frac{rS\,(K_{d2}+C+D)}{(Y-S)^2}
\ll1,
$$
by assumption 2.

\textbf{Step 6: $\theta_2\ll1$, $Y_f\ll K_{d2}$, and $\delta_2\ll1$.}
Because $\delta_2\ge0$, the factor $1/(1+\delta_2)$ in \eqref{eq:r3-3} is
at most 1, so
$$
\theta_2\le\frac{Y-S}{K_{d2}+C+D}\,(1+\delta_1)\ll1,
$$
since $(Y-S)/(K_{d2}+C+D)\ll1$ by assumption 3 and $1+\delta_1\approx1$ by
Step 5. To convert this into a statement about $Y_f$, start from
$Y_f=\theta_2(Y_f+K_{d2})$, subtract $\theta_2Y_f$ from both sides to get
$Y_f(1-\theta_2)=\theta_2K_{d2}$, and divide by $K_{d2}(1-\theta_2)$:
$$
\frac{Y_f}{K_{d2}}=\frac{\theta_2}{1-\theta_2}.
$$
With $\theta_2\ll1$, the right-hand side is approximately $\theta_2\ll1$,
so $Y_f\ll K_{d2}$. Finally, because $K_{d2}\le K_{d2}+C+D$,
$$
\delta_2=\frac{Y_f}{K_{d2}+C+D}\le\frac{Y_f}{K_{d2}}\ll1.
$$

\textbf{Step 7: result.} Since $\delta_1\ll1$ and $\delta_2\ll1$,
$(1+\delta_1)/(1+\delta_2)\approx1$, and \eqref{eq:r3-3} becomes
$$
\theta_2\approx\frac{Y-S}{K_{d2}+C+D}.
$$
Substituting this into \eqref{eq:r3-1}, and using $1-r\approx1$ (because
$K_{d1}\ll K_{d2}$),
\begin{equation}
\rho\approx r+\frac{Y-S}{K_{d2}+C+D}.
\label{eq:r3-5}
\end{equation}
Correspondingly, since $Y_f/K_{d2}=\theta_2/(1-\theta_2)\approx\theta_2$,
the free antibody is $Y_f\approx K_{d2}(Y-S)/(K_{d2}+C+D)$, confirming the
heuristic estimate. Substituting $S=A+B$ for the control and $S=A'+B$
for the treatment,
$$
\rho\approx r+\frac{Y-A-B}{K_{d2}+C+D},
\qquad
\rho'\approx r+\frac{Y-A'-B}{K_{d2}+C+D}.
$$

(An exact lower bound also follows. Since
$d\theta_2/dY_f=K_{d2}/(Y_f+K_{d2})^2>0$, $\theta_2$ increases with
$Y_f$, so substituting the lower bound \eqref{eq:r3-4} into
$\theta_2=Y_f/(Y_f+K_{d2})$ and then multiplying the numerator and
denominator by $(K_{d2}+C+D)/K_{d2}$,
$$
\theta_2
\ge\frac{K_{d2}\frac{Y-S}{K_{d2}+C+D}}
{K_{d2}\frac{Y-S}{K_{d2}+C+D}+K_{d2}}
=\frac{Y-S}{(Y-S)+(K_{d2}+C+D)}.
$$
Substituting into \eqref{eq:r3-1} gives
$\rho\ge r+(1-r)(Y-S)/(Y-S+K_{d2}+C+D)$, which agrees with
\eqref{eq:r3-5} to first order when $Y-S\ll K_{d2}+C+D$.)

Two terms contribute to $\rho$. The first, $r$, is the baseline ratio of
off-target to target capture efficiency set by the affinities. The
second is the fraction of off-target epitopes occupied by the excess
antibody. Because target is saturated ($\theta_1\approx1$), any increase
in off-target occupancy translates directly into an increase in $\rho$.
Unlike in Regime 1, the second term is not multiplied by $r$, so $\rho$ is
not necessarily close to $r$: it can be many times larger.

\textbf{Two sub-cases.} The denominator $K_{d2}+C+D$ reflects the two
places where the excess antibody can go, in the ratio $K_{d2}:(C+D)$.
\begin{itemize}
    \item[(3a)] Few off-target epitopes, $C+D\ll K_{d2}$. Then
    $K_{d2}+C+D\approx K_{d2}$, so
    $$
    Y_f\approx K_{d2}\frac{Y-S}{K_{d2}}=Y-S,
    \qquad
    \rho\approx r+\frac{Y-S}{K_{d2}}.
    $$
    The excess antibody stays free in solution, and the off-target
    occupancy is simply $Y_f/K_{d2}$. Assumption 3 becomes $Y-S\ll K_{d2}$,
    and assumption 2 becomes $K_{d1}S\ll(Y-S)^2$, the exact mirror of
    Regime 1's assumption 2.

    \item[(3b)] Many off-target epitopes, $C+D\gg K_{d2}$. Then
    $K_{d2}+C+D\approx C+D$, so
    $$
    Y_f\approx K_{d2}\frac{Y-S}{C+D}\ll Y-S,
    \qquad
    \rho\approx r+\frac{Y-S}{C+D}.
    $$
    The off-target pool soaks up nearly all of the excess antibody, so the
    excess divided by the number of off-target epitopes is the fraction of
    off-target epitopes that are occupied. Assumption 3 becomes
    $Y-S\ll C+D$.
\end{itemize}
Chromatin contains an enormous number of weak, nonspecific binding sites,
so sub-case (3b) is plausible in practice.

(Remark on assumption 3. The approximation to $\theta_2$ in
\eqref{eq:r3-3} actually requires only $\delta_2\ll1$, i.e.,
$Y_f\ll K_{d2}+C+D$. In sub-case (3b), this can remain true even when a
substantial fraction of off-target epitopes is occupied, because the
large off-target pool keeps $Y_f$ small. Then the exact relation
\eqref{eq:r3-1} with $\theta_2\approx(Y-S)/(C+D)$ still holds, and $\rho$
continues to grow linearly with the excess until $Y-S$ approaches $C+D$.
We keep assumption 3 as stated because it defines the regime in which
$\rho\ll1$.)

\textbf{How $\rho'$ compares to $\rho$.} Subtracting the approximations
for $\rho$ and $\rho'$, the $r$ terms cancel and the fractions share a
common denominator:
$$\begin{aligned}
\rho'-\rho
&\approx\left(r+\frac{Y-S_t}{K_{d2}+C+D}\right)
-\left(r+\frac{Y-S_c}{K_{d2}+C+D}\right) \\
&=\frac{(Y-S_t)-(Y-S_c)}{K_{d2}+C+D}
=\frac{S_c-S_t}{K_{d2}+C+D}.
\end{aligned}$$
Substituting $S_c=A+B$ and $S_t=A'+B$, the spike-in terms $B$ cancel:
\begin{equation}
\rho'-\rho\approx\frac{(A+B)-(A'+B)}{K_{d2}+C+D}
=\frac{A-A'}{K_{d2}+C+D}.
\label{eq:r3-6}
\end{equation}
Since the denominator is positive, \textbf{$\rho'>\rho$ when $A'<A$}.
Intuitively, the treatment removes $A-A'$ target epitopes, which frees
the $A-A'$ antibody molecules that would have occupied them. That
antibody joins the excess, and a fraction $(C+D)/(K_{d2}+C+D)$ of it
lands on off-target epitopes, raising the off-target occupancy.

The sign of $\rho'-\rho$ is in fact exact. Define
$g(Y_f)=Y_f+SY_f/(Y_f+K_{d1})+(C+D)Y_f/(Y_f+K_{d2})$, so that antibody
conservation reads $g(Y_f)=Y$. Each term of $g$ increases with $Y_f$,
and for $Y_f>0$ the middle term increases with $S$. If $S$ decreases
(from $S_c$ to $S_t$) while $Y$ is held fixed, then $g$ evaluated at the
old $Y_f$ falls below $Y$, so the new solution must have a larger $Y_f$.
By the quotient rule,
$$
\frac{d\rho}{dY_f}
=\frac{(Y_f+K_{d2})-(Y_f+K_{d1})}{(Y_f+K_{d2})^2}
=\frac{K_{d2}-K_{d1}}{(Y_f+K_{d2})^2}>0,
$$
so a larger $Y_f$ means a larger $\rho$. This argument applies in every
regime, including Regime 1.

Note that the magnitude \eqref{eq:r3-6}, unlike the Regime 1 result, is
not suppressed by a factor of $r$. In Regime 1, a decrease in target
changes $\rho$ only by changing how tightly the (limiting) antibody is
held on target, an effect proportional to $r$. Here, the decrease in
target releases antibody directly into the off-target pool.

\textbf{Relevance to practice.} Typical ChIP protocols use micrograms of
antibody. An IgG has a molar mass of about 150~kDa, so
$1~\mu\mathrm{g}$ is about
$10^{-6}~\mathrm{g}/(1.5\times10^{5}~\mathrm{g/mol})\approx6.7~\mathrm{pmol}$.
By comparison, a factor with $10^4$ binding sites per cell in $10^6$ cells
has $10^{10}$ sites, or about
$10^{10}/(6.0\times10^{23}~\mathrm{mol}^{-1})\approx0.017~\mathrm{pmol}$.
For many targets, then, standard practice places the IP in this
antibody-excess regime. (Very abundant epitopes, such as some histone
modifications, may instead be comparable to or exceed the antibody.)

\subsubsection{ChIP normalization}

The read-count model of Regime 1 applies unchanged, because it used only
the composition $A:B:\rho C:\rho D$ of the immunoprecipitation. In
particular, ChIP-Rx scaling gives
$$
\frac{n_i'/N_s'}{n_i/N_s}
=\frac{a_i'}{a_i}\cdot\frac{B+\rho D}{B+\rho' D}.
$$

\textbf{Rewriting the bias factor.} Write $\rho'=\rho+(\rho'-\rho)$, so
the denominator is
$B+\rho'D=(B+\rho D)+(\rho'-\rho)D$. Dividing the numerator and
denominator of the bias factor by $B+\rho D$,
$$
\frac{B+\rho D}{B+\rho'D}
=\frac{B+\rho D}{(B+\rho D)+(\rho'-\rho)D}
=\frac{1}{1+\dfrac{(\rho'-\rho)D}{B+\rho D}}.
$$
Since $\rho'>\rho$ when $A'<A$, the fraction in the denominator is
positive, so the bias factor is less than 1: as in Regime 1,
\textbf{ChIP-Rx overstates the decrease} in the treatment condition. The
bias factor is approximately 1 only if
$$
(\rho'-\rho)D\ll B+\rho D.
$$

\textbf{Substituting the regime expressions.} Substituting
\eqref{eq:r3-6} for $\rho'-\rho$ and \eqref{eq:r3-5} for $\rho$, the
bias factor is
$$
\frac{B+\rho D}{B+\rho'D}
\approx\frac{1}{1+\dfrac{(A-A')D}{(K_{d2}+C+D)\bigl(B+\rho D\bigr)}},
\qquad
\rho\approx r+\frac{Y-A-B}{K_{d2}+C+D},
$$
and ChIP-Rx is unbiased only if
$$
\frac{(A-A')D}{K_{d2}+C+D}\ll B+\left(r+\frac{Y-A-B}{K_{d2}+C+D}\right)D.
$$
The left-hand side has a direct interpretation. Splitting the fraction,
$$
\frac{(A-A')D}{K_{d2}+C+D}
=(A-A')\cdot\frac{C+D}{K_{d2}+C+D}\cdot\frac{D}{C+D}:
$$
the $A-A'$ antibody molecules freed by the treatment, times the fraction
of excess antibody that is captured off-target, times the share of
off-target epitopes that belong to the spike-in. That is the extra
antibody that ends up on spike-in off-target epitopes in the treatment.
The right-hand side, $B+\rho D=(B\theta_1+D\theta_2)/\theta_1$, is the
spike-in material in the control IP divided by $\theta_1\approx1$. So the
condition says that the extra spike-in capture caused by the freed
antibody must be small compared with the spike-in capture already
present.

\textbf{How large can the bias be?} Consider sub-case (3b), so
$K_{d2}+C+D\approx C+D$, and suppose off-target spike-in reads are a
minor part of the spike-in signal, $\rho D\ll B$. The fraction in the
bias factor becomes
$$
\frac{(A-A')D}{(K_{d2}+C+D)(B+\rho D)}
\approx\frac{(A-A')D}{(C+D)B}.
$$
Suppose further that the spike-in chromatin has the same ratio of
off-target to target epitopes as the cells, $D/B=C/A$. Substituting
$D/B=C/A$,
$$
\frac{(A-A')D}{(C+D)B}
=\frac{(A-A')}{C+D}\cdot\frac{C}{A}
=\frac{A-A'}{A}\cdot\frac{C}{C+D}.
$$
If the spike-in is a small fraction of the chromatin ($D\ll C$), then
$C/(C+D)\approx1$ and the fraction is approximately $1-A'/A$, the
fractional loss of target. The bias factor is then
$$
\frac{1}{1+(1-A'/A)}=\frac{1}{2-A'/A},
$$
and the measured global fold change is
$$
\frac{A'}{A}\cdot\frac{1}{2-A'/A}.
$$
For example, if the treatment halves the target ($A'/A=1/2$), the bias
factor is $1/(2-1/2)=2/3$, and ChIP-Rx reports a fold change of
$(1/2)(2/3)=1/3$ instead of $1/2$. This bias is of order 1, whereas in
Regime 1 it was suppressed by $r$. The reason is that in this regime the
freed antibody is redistributed almost entirely onto off-target
epitopes, and the spike-in receives a share of that redistribution
comparable to its share of the target signal. (If instead $\rho D\gg B$,
the bias factor approaches $\rho D/(\rho' D)=\rho/\rho'$, the ratio of
off-target capture efficiencies, which can also differ substantially
from 1 when the excess $Y-S$ is not large compared with $A-A'$.)

\textbf{Scaling by spike-in target reads.} The derivation of Regime 1
used only the read-count model, not any regime-specific approximation.
Therefore, counting only reads $N_{s,B}$ at the spike-in's target loci
still gives
$$
\frac{n_i'/N_{s,B}'}{n_i/N_{s,B}}=\frac{a_i'}{a_i},
$$
removing the bias regardless of $\rho$ and $\rho'$. This matters more
here than in Regime 1, because the off-target spike-in reads that
$N_s$ includes are exactly what changes between conditions. (As before,
reads in spike-in peaks include some off-target background, so the
cancellation is only approximate; the residual error is larger when
$\rho D$ is larger.)
